\documentclass[10pt,a4paper]{article}

\usepackage[margin=20mm]{geometry}
\usepackage{amsmath,amssymb,amsthm}

\theoremstyle{definition}
\newtheorem*{problem*}{Problem}

\setlength{\parindent}{0pt}
\setlength{\parskip}{5pt}

\begin{document}

\begin{center}
  {\Large\bfseries Degenerate Cahn--Hilliard Convergence to Surface Diffusion}
\end{center}

\section{Definitions}

Let $\mathbb T:=\mathbb R/(2\pi\mathbb Z)$.

Let
\begin{equation}
  \Omega=\mathbb T^2,
\end{equation}
and let $E_0\Subset\Omega$ have smooth embedded boundary
\begin{equation}
  \Gamma_0=\partial E_0.
\end{equation}
Let
\begin{equation}
  \Gamma(t)=\partial E(t),
  \qquad
  0\leq t<T_*,
\end{equation}
be the maximal smooth surface-diffusion flow
\begin{equation}
  V=-\Delta_\Gamma H,
  \qquad
  H=-\operatorname{div}_\Gamma\nu,
\end{equation}
where $\nu$ is outward. Fix $T<T_*$. Define
\begin{equation}
  F(u)=\frac{1}{2}(1-u^2)^2,
  \qquad
  q(z)=\tanh z,
  \qquad
  m_\varepsilon(u)=(1-u^2)^2+\varepsilon^6.
\end{equation}
On the slow surface-diffusion time scale, consider
\begin{equation}\label{eq:ch}
  \varepsilon^2\partial_tu_\varepsilon
  =
  \frac{9}{4}
  \operatorname{div}
  \bigl(m_\varepsilon(u_\varepsilon)\nabla\mu_\varepsilon\bigr),
  \qquad
  \mu_\varepsilon
  =
  -\varepsilon^2\Delta u_\varepsilon+F'(u_\varepsilon).
\end{equation}
Let $d_0$ be a smooth truncation of the signed distance to $\Gamma_0$,
positive in $E_0$. Let $H_0^e$ be a smooth normal extension of $H_0$ in a
fixed tubular neighborhood, cut off outside that neighborhood. Take
\begin{equation}
  u_\varepsilon^0(x)
  =
  q\left(\frac{d_0(x)-a_\varepsilon}{\varepsilon}\right)
  -\frac{\varepsilon}{6}H_0^e(x),
\end{equation}
where the unique small shift $a_\varepsilon$ is chosen so that
\begin{equation}
  \int_\Omega u_\varepsilon^0\,dx
  =
  2|E_0|-|\Omega|.
\end{equation}

\section{Best Known Result}

The matched expansion of the well-prepared initial layer gives
\begin{equation}
  \mu_\varepsilon^0
  =
  -\frac{2}{3}\varepsilon H_0+O(\varepsilon^2)
\end{equation}
inside the interfacial region. Together with the factor $9/4$, this formally
selects the surface-diffusion law $V=-\Delta_\Gamma H$. For every classical
phase-field solution, the normalized energy obeys the exact identity
\begin{equation}
\begin{split}
  \frac{d}{dt}
  \left[
    \frac{3}{4}\int_\Omega
    \left(
      \frac{\varepsilon}{2}|\nabla u_\varepsilon|^2
      +\frac{1}{\varepsilon}F(u_\varepsilon)
    \right)dx
  \right]
  ={}&
  -\frac{27}{16\varepsilon^3}
  \int_\Omega
  m_\varepsilon(u_\varepsilon)|\nabla\mu_\varepsilon|^2\,dx.
\end{split}
\end{equation}
A flat transition layer already has an $L^1$ error of order
$2(\log2)\varepsilon$ per unit interface length, so an $O(\varepsilon)$
phase-field estimate has the optimal natural order.

\newpage
\section{Problem Formalization}

\begin{problem*}[Degenerate Cahn--Hilliard limit]
Prove that, for every $T<T_*$ and all sufficiently small $\varepsilon$,
equation~\eqref{eq:ch} has a classical solution on $[0,T]$ and that its zero
level set
\begin{equation}
  \Gamma_\varepsilon(t)
  =
  \{u_\varepsilon(t)=0\}
\end{equation}
is a single embedded normal graph over $\Gamma(t)$.

Prove the optimal-order phase-field estimate
\begin{equation}
  \sup_{0\leq t\leq T}
  \left\|
    u_\varepsilon(t)
    -
    \bigl(2\chi_{E(t)}-1\bigr)
  \right\|_{L^1(\Omega)}
  \leq C_T\varepsilon,
\end{equation}
and the Hausdorff convergence
\begin{equation}
  \sup_{0\leq t\leq T}
  d_H\bigl(\Gamma_\varepsilon(t),\Gamma(t)\bigr)
  \longrightarrow 0.
\end{equation}

Define the normalized energy measures
\begin{equation}
  \lambda_\varepsilon^t
  =
  \frac{3}{4}
  \left(
    \frac{\varepsilon}{2}|\nabla u_\varepsilon(t)|^2
    +
    \frac{1}{\varepsilon}F(u_\varepsilon(t))
  \right)dx.
\end{equation}
Prove that
\begin{equation}
  \lambda_\varepsilon^t
  \stackrel{*}{\rightharpoonup}
  \mathcal H^1\!\llcorner\Gamma(t)
\end{equation}
uniformly for $t\in[0,T]$ against every fixed continuous test function.

Finally, prove convergence of the total dissipation:
\begin{equation}
  \frac{27}{16\varepsilon^3}
  \int_0^T\int_\Omega
  m_\varepsilon(u_\varepsilon)|\nabla\mu_\varepsilon|^2\,dx\,dt
  \longrightarrow
  \int_0^T\int_{\Gamma(t)}
  |\nabla_\Gamma H|^2\,ds\,dt.
\end{equation}
\end{problem*}

\end{document}
